\section{Status-preserving budgets and the final contradiction}\label{sec:finish}
We spell out the probability-kernel replacements and the precise deterministic
inputs reused from S Sections 3--8. Their charts, label budget, finite-tree
count, posterior selection argument, shared-union detour and final filters
are antecedents, not new claims. All events below live on $(t,V)$ and use
cap distances and cap labels. Remark~\ref{rem:kernels} supplies reference
and label stability for the finite kernels used here on their upward events.

\subsection{Charts, comparison, and the one-edge budget}
S Section 3's deterministic chart construction gives countably many compact
label charts $I$, each with a horizontal order parameter $r(\rho)$, finite
width $W_I=r(\rho_r)-r(\rho_\ell)$, and slope $\sigma>0$. The arm cones are
\[
 Q^\pm(k,H)=\{z:\ \pm(z_2-k_2)\ge\sigma|z_1-k_1|-H\}.
\]
Choose nested finite deterministic endpoint sets $D_j$ dense in chart
order and containing the endpoints of every complementary gap. Consecutive
endpoints define the bins $\Pi_j$. Their widths sum to $W_I$, and distinct
nonendpoint labels eventually separate. Fixed-label backwards finiteness
excludes bigeodesic end labels at the countable union of those endpoints.

Here is the deterministic comparison from S Section 5. Suppose a finite
increase near a site $k$ adds total cost at least $b$ on finite initial
portions of a bigeodesic whose upper label lies in a bin $(\alpha,\beta)$.
Suppose its arms have cone offset $H_p$, and the base upper $\alpha$ and
lower $-\beta$ comparator rays from $y$ have offset $H_c$. If the change
and gain portions lie in the radius-$R$ box about $k$, then for
\[
 y_1-k_1>S=\max\left\{R+(R+H_c)/\sigma,(H_p+H_c)/\sigma\right\}
\]
the comparator rays avoid the changed box, remain geodesics with their
labels, and the chart crossing argument gives
\begin{equation}\label{eq:budget}
 b\le B_\alpha(y,k)+B_{-\beta}(y,k).
\end{equation}
This uses finite uniqueness, cone crossing and calibration. Successive
comparisons telescope, leaving only the horizontal label width. The cap
has every one of these premises by Sections~\ref{sec:labels}--\ref{sec:splice}.

At a grid edge $e$, independently replace the original $t_e$ by a sample
from its status-conditional law. Each individual output has exactly the
original iid law, with the same independent $V$. All statuses and thresholds
stay fixed, so only that cap edge changes. Conditional on the \emph{entire}
individual output and $V$, its hidden old value is in the low band of its
status $i$ with probability $p_{a,i}/r_i\ge p_*$. On a usable output
bigeodesic the edge is wholly uncapped by Proposition~\ref{prop:equivalence}.
Mark when its hidden old value is low. The old value is below $a_i<b_i$
and the common threshold exceeds the output value, so the cap increase
is at least $g$ and no other cap changes.

The source's base-comparator regularity check uses both endpoints of $e$:
a base ray either avoids it or uses it first, in which case its tail starts
at the other endpoint. Monotonicity and label transfer give the required
base cone offset from output regularity. Thus its spacing
$\ell(H)\le C_0(H+1)$ and \eqref{eq:budget} apply. If
$A_e^{\alpha,\beta}(H)$ denotes a usable, $H$-regular output bigeodesic
event with label in that bin, the mark probability and budget imply
\begin{equation}\label{eq:edgecount}
 \PP(A_e^{\alpha,\beta}(H))
 \le\frac{\ell(H)\,[r(\beta)-r(\alpha)]}{g p_*},
 \qquad \EE N_e(H)\le C(H+1).
\end{equation}
The latter follows by partition refinement and monotone convergence; $N_e(H)$
counts distinct usable regular edge labels. Independence of marks is not needed.

S Section 6's rooted finite-tree leaf count credits at least $nF_n/4$
distinct usable edge-label pairs to $F_n$ forks. Its usable-path estimate
is valid here in the \emph{original} iid field: a path of length $j$ has
at least half bad edges with probability at most $2^j\epsilon^{j/2}$,
and the path count contributes at most $4^j$, giving $(8\sqrt\epsilon)^j$.
Cap bigeodesics are wholly unchanged. The simultaneous cap cone bound
and \eqref{eq:edgecount} therefore provide exactly its expectation bound.
First let $n\to\infty$ at fixed cone error and then let that error vanish.
There are no forks. With opposite exceptional coalescence from
Section~\ref{sec:splice}, the source's deterministic conclusion identifies
the same two Busemann functions for all bigeodesics of a common random
label, simultaneously. This is not an intersection of fixed-label null events.

\subsection{Finite raises, likelihoods, and posterior independence}
Assume bigeodesics have positive probability. The source's countable chart
and cone-offset reduction gives fixed $I,\sigma,H_0$ and $p>0$ for a
bigeodesic through the origin with arms in $Q^\pm(0,H_0)$ and nonendpoint
label in $I$. Shrink this event, if necessary, by fixing a usable-prefix
threshold $n_0$; the path count above makes the discarded probability
arbitrarily small. Retain the notation $p$ for the resulting positive
probability. These choices are made before the requested toll $m$.

Let $h(e)=\max(|w_2|,|w'_2|)$ for the endpoints of $e$, and put
\[
 \delta_h=\frac{\eta}{(h+2)\log(h+2)},\qquad
 A_k(R)=\left\{e:h(e)\le R,\ 
 \max_{w\in e}|w_1-k_1|\le1+\frac{h(e)+H_0}{\sigma}\right\}.
\]
Take independent trigger uniforms and samples $Z_{e,i}\sim G(\cdot\mid H_i)$,
independent of $X$ and across edges. If $t_e\in L_i$, raise it to $Z_{e,i}$
with probability $\delta_h/p_{a,i}$; otherwise keep it. Choose $\eta$ so
$\delta_0\le\min_i p_{a,i}/2$. Use the same sample for every region
containing $e$, applying it only once in overlaps. All statuses and all
$V$ remain fixed, so all $K_e$ remain fixed. Cap changes are confined
exactly to the raised edges of the specified finite region.

The one-edge output density relative to $G$ is
\begin{align}
 \lambda_h(t)&=1+\sum_{i=o,c}\left[-\frac{\delta_h}{p_{a,i}}\ind_{L_i}(t)
                     +\frac{\delta_h}{p_{b,i}}\ind_{H_i}(t)\right],\label{eq:density}\\
 \int\lambda_h\,dG&=1,\qquad
 \int\lambda_h^2\,dG=1+\delta_h^2\sum_{i=o,c}
                       (p_{a,i}^{-1}+p_{b,i}^{-1}).\nonumber
\end{align}
There is mass $\delta_h$ moved in each status, not just in one status.
The number of region edges at height $h$ is at most
$(12+8/\sigma)(h+H_0+1)$. Independence gives a product likelihood $L$ with
\[
 \EE L^2\le\exp(C\eta^2),\qquad
 \TV(\mathcal L(Y^{k,R}),\mathcal L(X))
 \le\tfrac12\sqrt{e^{C\eta^2}-1},
\]
uniformly in $k,R$, because
$\sum_h(h+H_0+1)/((h+2)^2\log^2(h+2))<\infty$.
Choose $\eta$ so this is at most $p/8$. Finite union fields are also
absolutely continuous. Cap pushforward cannot increase total variation.

Conditional on the whole original output $Y$ and all cap auxiliaries $V$,
the hidden success indicators are independent, with
\begin{equation}\label{eq:posterior}
 \theta_e=\sum_{i=o,c}\ind_{H_i}(Y_e)
                  \frac{\delta_{h(e)}}{p_{b,i}+\delta_{h(e)}}.
\end{equation}
Indeed the joint success/output measure in $H_i$ is
$(\delta_h/p_{b,i})G(dt)$; divide by \eqref{eq:density}. The product
conditional law extends from finite output cylinders to the full output
sigma-field. Circuit lengths are already functions of the output statuses,
so they introduce no additional conditioning. A witnessing cap bigeodesic
is selected from $(Y,V)$ alone. On a successful usable edge of it,
$Y_e=s_e\le K_e$ and the old original is at most $a_i<b_i\le Y_e$;
the old cap is the old original too. Thus the cap gain is at least $g$.

For every prefix length at least $n_0$, at least half its edges are usable.
The $j$th upper-arm edge has $h(e)\le j$, so Abel summation and
\eqref{eq:posterior} give a conditional expected number of successes
bounded below by a positive constant times
$\sum_{j=n_0}^R((j+2)\log(j+2))^{-1}$, which diverges. Conditional
variance is at most the mean. Fix a near radius $R_0$ with mean at least
eight. For any toll $m$, choose $R(m)$ with mean at least
$\max(2\lceil m/g\rceil,32)$. Chebyshev bounds give, jointly, a near
strict raise and total gain at least $m$ with conditional probability at
least $3/4$. The source's marking calculation then gives probability at
least $p/4$. The horizontal spacing
\[
 \ell>R_0+1+(R_0+H_0)/\sigma
\]
is fixed independently of $m$. Posterior independence is not asserted
after the later base or neighbor filters.

\subsection{Shared unions and telescoping}
A bigeodesic in the individual region at $k$ has every edge of height
$\le R$ inside $A_k(R)$ by its cone width; higher edges are never raised.
Passing to the two-site union therefore leaves all of its edges unchanged,
while competitors only increase. The same holds for the selected path at
the second site. A near raised edge of the second path lies outside the
first region by the fixed spacing, so in the first field it still has its
base cap cost and in the union its cost is strictly larger. The common
Busemann pair for equal labels and the source's vanishing-excess detour
through this strict edge rule out repeated marked labels. This is exactly
the geometric use of the near raise; the large radius pays only the toll.

S Section 8's base-comparator filter loses at most $p/8$, leaving
$q_m\ge p/8$. For each fixed $m$, the comparison separation $S(m)$ is finite.
Refine $\Pi_j$ and discard marks with a same-bin neighbor within $S(m)$.
Distinct labels and finitely many neighbors make this filter disappear
pointwise as $j\to\infty$. Apply \eqref{eq:budget} successively in each
bin; Busemann increments telescope and their linear parts leave the sum
$W_I$ of bin widths. The bound is $q_{m,j}\le\ell W_I/m$. Bounded convergence
therefore yields
\[
             p/8\le q_m\le\ell W_I/m\qquad\text{for every integer }m.
\]
The constants $p,\ell,W_I$ were fixed before $m$, a contradiction for large
$m$. Sites, finite regions, unions, partition endpoints and chosen radii
form countable collections, so all inherited full events can be imposed
simultaneously. Random-label equalities were already pathwise.
There are no cap bigeodesics. Proposition~\ref{prop:equivalence} proves
Theorem~\ref{thm:main} for the remaining laws, and S covers \eqref{eq:paid}.
